\documentclass[11pt]{article}

\usepackage[margin=1in]{geometry}
\usepackage[T1]{fontenc}
\usepackage[utf8]{inputenc}
\usepackage{lmodern}
\usepackage{amsmath,amssymb,amsfonts}
\usepackage{xcolor}
\usepackage[hidelinks]{hyperref}
\usepackage{booktabs}
\usepackage{enumitem}
\usepackage{titlesec}

\setlength{\parskip}{2pt}
\titlespacing*{\section}{0pt}{8pt plus 2pt minus 2pt}{4pt}
\titlespacing*{\paragraph}{0pt}{4pt plus 1pt minus 1pt}{6pt}

\newcommand{\R}{\mathbb{R}}
\newcommand{\E}{\mathbb{E}}
\newcommand{\Cov}{\mathrm{Cov}}
\newcommand{\sign}{\mathrm{sign}}
\newcommand{\Tsel}{T_{\mathrm{sel}}}
\newcommand{\Tscore}{T_{\mathrm{score}}}
\newcommand{\dout}{d_{\mathrm{out}}}

\begin{document}

\title{\Large\textbf{Soft LUT gradients in dominance space}\\[2pt]
\large Math notes toward a ranking-aligned weight update}
\author{Anatoly Starostin}
\date{}
\maketitle
\vspace{-14pt}

% =========================================================================
\section{Setup}
\label{sec:setup}

Single \texttt{TinyMHLut} table (drop head / tph indices). Input $x \in \R^E$, $n_c$ anchor pairs $(i_j, k_j)$ for $j = 1, \dots, n_c$, output $y \in \R^{\dout}$. Pair differences gathered into the linear operator
\[
D : \R^E \to \R^{n_c}, \qquad (Dx)_j = x_{i_j} - x_{k_j} =: z_j, \qquad \text{rows of } D = e_{i_j} - e_{k_j}.
\]
Pair scores $u_j = z_j / \Tscore$. Table $T \in \R^{K \times \dout}$ with $K = 2^{n_c}$ rows indexed by $b \in \{0,1\}^{n_c}$. Write the parity sign
\[
\chi_j(b) := 2 b_j - 1 \in \{-1, +1\}, \qquad \chi_S(b) := \prod_{j \in S} \chi_j(b).
\]

% =========================================================================
\section{Forward: softmax over soft signs}
\label{sec:fwd}

Row energy and selection distribution:
\[
\mathcal{E}(b) = \frac{1}{\Tsel} \sum_j u_j\, \chi_j(b), \qquad
p(b) = \frac{e^{\mathcal{E}(b)}}{Z}, \qquad
Z = \sum_b e^{\mathcal{E}(b)}.
\]
Output:
\[
y \;=\; \E_p\!\left[T_b\right] \;=\; \sum_b p(b)\, T_b.
\]

\paragraph{Factorisation.} Because $\mathcal{E}$ is a sum over independent $\pm$ terms,
\[
Z \;=\; \prod_j 2 \cosh\!\big(u_j/\Tsel\big), \qquad
p(b) \;=\; \prod_j q_j^{b_j}\,(1 - q_j)^{1-b_j}, \qquad
q_j \;=\; \sigma\!\big(2 u_j / \Tsel\big).
\]
The $K$-way softmax is equivalent to an $n_c$-fold product of independent Bernoullis. We never need to materialise $p \in \R^K$ — only $\{q_j\} \in [0,1]^{n_c}$.

% =========================================================================
\section{Standard backward}
\label{sec:bwd}

Let $g := \partial L / \partial y \in \R^{\dout}$.

\paragraph{Weight gradient ($\partial L / \partial T$).}
\[
\frac{\partial L}{\partial T_b} \;=\; p(b)\, g.
\]
Pointwise on rows, weighted by selection probability. \emph{Dense in support, sparse in mass:} every row receives a gradient, but most $p(b)$ are vanishingly small, so the useful signal lives on the few high-probability rows.

\paragraph{Input gradient ($\partial L / \partial x$).} By the chain rule,
\[
\frac{\partial L}{\partial x} \;=\; D^\top \frac{\partial L}{\partial z}, \qquad
\frac{\partial L}{\partial z_j} \;=\; \frac{1}{\Tscore} \frac{\partial L}{\partial u_j} \;=\; \frac{1}{\Tscore}\, g^\top \frac{\partial y}{\partial u_j}.
\]
Differentiating $p$ in the softmax form gives
\[
\frac{\partial p(b)}{\partial u_j} \;=\; \frac{p(b)}{\Tsel}\,\big(\chi_j(b) - \E_p[\chi_j]\big),
\]
so
\[
\frac{\partial y}{\partial u_j} \;=\; \frac{1}{\Tsel}\, \Cov_p\!\big(T_b,\, \chi_j(b)\big).
\]
Because $\chi_j$ depends on the single bit $b_j$ (marginal $q_j$ under $p$), the $K$-sum collapses to a difference of two conditional row-means:
\[
\boxed{\;\Cov_p(T_b, \chi_j) \;=\; 2\, q_j(1 - q_j)\, \Big(\E_{p \mid b_j = 1}[T_b] \;-\; \E_{p \mid b_j = 0}[T_b]\Big).\;}
\]
Combining:
\[
\frac{\partial L}{\partial u_j} \;=\; \frac{2\, q_j(1-q_j)}{\Tsel}\, g^\top\!\big(\E_{p \mid b_j=1}[T_b] - \E_{p \mid b_j=0}[T_b]\big).
\]

\paragraph{Geometric note.} $\partial L / \partial x$ lives in the row space of $D$ --- a subspace of $\R^E$ of dimension at most $n_c$, the \emph{comparison-visible subspace}. Any input perturbation orthogonal to it is invisible to the LUT.

% =========================================================================
\section{Dominance-space reading: when does a ranking change?}
\label{sec:rank}

The $n_c$ hyperplanes $\{x : x_{i_j} = x_{k_j}\}$ partition $\R^E$ into chambers, each labelled by a sign pattern $b \in \{0,1\}^{n_c}$. The forward map collapses each chamber to its row $T_b$.

A step $\delta = -\eta\, \partial L / \partial x$ in the loss-decreasing direction flips pair $j$ when $z_j + (D \delta)_j$ crosses zero, i.e.\ at
\[
\eta^*_j \;=\; \frac{z_j}{(D\, \partial L / \partial x)_j} \;=\; \frac{z_j}{\partial L / \partial z_j}, \qquad \text{flip iff } \eta^*_j > 0.
\]
The first ranking change is reached at
\[
\eta^* \;=\; \min_{j\, :\, \eta^*_j > 0} \eta^*_j, \qquad j^* \;=\; \arg\min.
\]

\paragraph{Reading.}
\begin{itemize}
\setlength{\itemsep}{1pt}
    \item If $\sign(\partial L / \partial z_j) = -\sign(z_j)$: loss wants to widen the gap, \emph{strengthening} the current bit. This pair never flips ($\eta^*_j < 0$).
    \item Otherwise loss wants to flip the bit; the smallest $\eta^*_j$ among flip-wanting pairs identifies the very next ranking $b' = b \oplus e_{j^*}$.
\end{itemize}
This is the structural fact: $\partial L / \partial x$ encodes an \emph{itinerary through the chamber graph}, with $j^*$ being the next bit-flip the loss is reaching for.

% =========================================================================
\section{Two channels of $\partial L / \partial T$ --- and what's missing}
\label{sec:channels}

Differentiating $y = \E_p[T_b]$ splits the loss change into two channels:
\[
\delta y \;=\; \underbrace{\E_p\!\left[\delta T_b\right]}_{\text{(C1) change the table}} \;+\; \underbrace{\sum_b \big(\delta p(b)\big)\, T_b}_{\text{(C2) reroute by changing }x}.
\]
(C2) is what flows up through $\partial L / \partial x$ to the upstream layer. (C1) is the table's own update. \textbf{Standard backprop weights both channels by $p$ alone}, throwing away the ranking direction $j^*$ that channel (C2) just produced.

\emph{The idea.} Use $j^*$ to densify and structure the (C1) table update.

% =========================================================================
\section{Ranking-aligned weight update}
\label{sec:proposal}

Define a transport kernel $\Psi(b' \mid b, \partial L / \partial z)$ on the cube and set
\[
\widetilde{\frac{\partial L}{\partial T_{b'}}} \;=\; \Psi(b' \mid b,\, \partial L / \partial z)\, g.
\]
Three natural choices, ordered by structure:
\begin{enumerate}
\setlength{\itemsep}{2pt}
    \item \textbf{Hard STE.} $\Psi(b' \mid b) = \delta_{b, b'}$. Only the visited row updated.
    \item \textbf{Soft.} $\Psi(b' \mid b) = p(b')$. Every row touched, weighted by current confidence.
    \item \textbf{Ranking-aligned.} $\Psi(b' \mid b) = p(b') + \alpha\, \chi_{j^*}(b)\, \chi_{j^*}(b')$, or more generally
    \[
    \Psi(b' \mid b) \;=\; p(b')\, \exp\!\Big({-}\alpha \sum_{j\, :\, \eta^*_j > 0} \kappa_j\, \chi_j(b)\, \chi_j(b')\Big)
    \]
    with $\kappa_j$ decreasing in $\eta^*_j$ (closer-to-flipping pairs get more mass on the opposite bit).
\end{enumerate}

% =========================================================================
\section{The clean basis: Walsh characters on $\{0,1\}^{n_c}$}
\label{sec:walsh}

Every table is a linear combination of $K$ Walsh characters:
\[
T_b \;=\; \sum_{S \subseteq [n_c]} w_S\, \chi_S(b), \qquad
w_S \;=\; \frac{1}{K} \sum_b \chi_S(b)\, T_b,
\]
orthogonal in the cube measure: $\tfrac{1}{K}\sum_b \chi_S(b)\, \chi_{S'}(b) = \delta_{S, S'}$. A row-space update induces a Walsh-coefficient update
\[
\delta w_S \;=\; \frac{1}{K}\sum_b \chi_S(b)\, \delta T_b.
\]
The three kernels above produce:

\begin{center}
\renewcommand{\arraystretch}{1.2}
\begin{tabular}{lll}
\toprule
update style & $K \cdot \delta w_S \;\propto$ & character \\
\midrule
hard STE                  & $\chi_S(b)\, g$                       & broadband (all $w_S$ shifted equally) \\
soft backward             & $\E_p[\chi_S]\, g$                    & low-order $S$ dominate \\
ranking-aligned ($\alpha \to 1$)  & $\delta_{S,\{j^*\}}\, g$              & one-coefficient ($S = \{j^*\}$) \\
\bottomrule
\end{tabular}
\end{center}

Equivalently, translating the one-coefficient update back to row space:
\[
\boxed{\;\delta T_{b'} \;=\; -\eta\, g\, \chi_{j^*}(b') \;=\; -\eta\, g\,(2 b'_{j^*} - 1).\;}
\]
Every row's $j^*$ bit acts as the sign telling whether to add or subtract $\eta\, g$. The whole table is touched; the update is dense in rows but \emph{rank-one in Walsh basis}.

% =========================================================================
\section{What this corresponds to as a ``true'' gradient}
\label{sec:framings}

The above is not $\partial L / \partial T$ at the current routing $p$ --- it is a structured replacement. Three coherent ``it is the gradient of something'' framings:

\paragraph{(F1) Anticipatory / lookahead.}
It is the standard gradient of $L(\widehat y)$ where $\widehat y = \E_{\widehat p}[T_b]$ and $\widehat p$ is the predicted $k$-step-ahead routing distribution under SGD. The shift $p \to \widehat p$ deposits extra mass on the about-to-be-visited neighbour $b \oplus e_{j^*}$. Honest gradient descent on a slightly modified loss.

\paragraph{(F2) Natural gradient.}
Place a metric on table space
\[
\|\delta T\|_M^2 \;=\; \sum_S \lambda_S\, \|\delta w_S\|^2
\]
with $\lambda_S$ growing in $|S|$ (penalise high-frequency Walsh components). The natural gradient $-F^{-1} \nabla L$ then redistributes a pointwise row update along the low-frequency Walsh axes; the one-bit $\chi_{j^*}$ axis is the most natural single-coordinate move.

\paragraph{(F3) Variance reduction.}
The per-token soft gradient $p(b)\, g$ has large variance across tokens (one row dominates). Project onto $\chi_{j^*}$ as a control variate that captures the dominant low-order component; the residual is batch-averaged away.

All three roads produce the same shape of update; they justify it from different priors.

% =========================================================================
\section{Sanity checks and next steps}
\label{sec:next}

\begin{enumerate}
\setlength{\itemsep}{1pt}
    \item \textbf{Limiting behaviour.} When no clear $j^*$ exists (all $|\eta^*_j|$ comparable), the ranking-aligned kernel should collapse to standard soft backward. Verify on a 2-pair toy table.
    \item \textbf{Choice of $\Psi$.} One-bit ($\{j^*\}$ only), low-Hamming with decay, or full Walsh-soft
    $\Psi(b' \mid b) \propto p(b')\, e^{-\alpha d_H(b, b')}$.
    \item \textbf{Implementation.} A one-Walsh-axis update on $K$ rows is an outer product of two $\{-1, +1\}^{n_c}$ vectors; never need to materialise the full table.
    \item \textbf{Empirical probe.} On the smallest LUT in \texttt{exp428} (\texttt{qkv\_lut}, $\mathrm{NAP} = 6$, $\mathrm{tph} = 64$), swap soft backward for the $\chi_{j^*}$-aligned variant and compare bpb trajectories at $b = 16$.
\end{enumerate}

\end{document}
